\documentclass[11pt,a4paper]{article}

% ── Encoding & language ──
\usepackage[utf8]{inputenc}
\usepackage[T1]{fontenc}
\usepackage[english]{babel}

% ── Math & symbols ──
\usepackage{amsmath,amssymb,amsthm}

% ── Algorithms ──
\usepackage[ruled,vlined,linesnumbered]{algorithm2e}

% ── Tables ──
\usepackage{booktabs}
\usepackage{array}

% ── Links ──
\usepackage[colorlinks=true,linkcolor=blue,citecolor=blue,urlcolor=blue]{hyperref}
\usepackage{url}

% ── Layout ──
\usepackage[margin=2.5cm]{geometry}

% ── Listings for pseudocode ──
\usepackage{listings}
\lstset{
  basicstyle=\ttfamily\small,
  breaklines=true,
  frame=single,
  xleftmargin=1em
}

% ── Theorems / definitions ──
\newtheorem{definition}{Definition}
\newtheorem{thesis}{Design Thesis}
\newtheorem{condition}{Condition}

% ── Bibliography ──
\usepackage[numbers,sort&compress]{natbib}

% ══════════════════════════════════════════
\title{Controlled Generative Development:\\
A Deterministic Structural Core\\
for AI-Native Software Design}

\author{Alexander Fedorenko\\
\textit{Independent Researcher}\\
\texttt{alex.grig.fed@gmail.com}}

\date{March 2026}

\begin{document}
\maketitle

% ── Abstract ──
\begin{abstract}
This paper proposes Controlled Generative Development (CGD), a contract-first
framework for AI-assisted development in which the specification precedes code,
governs generation, and outlives code as the canonical artifact.  The central
object of CGD is the formal triple $T = \langle Y, G, \Sigma \rangle$, where
$Y$ is expressed through YASF, a minimal authoring layer in controlled natural
language, $G$ is a typed directed graph of functions and tables, and $\Sigma$
is a set of table signatures.  The paper formalizes the six-component function
contract, subgraph contract computation, decomposition and aggregation rules,
task-driven context assembly, and verification-driven dialogue.  CGD is
positioned as a contract-first approach between prompt-first, workflow-first,
spec-first, and formal-first traditions.  Using an end-to-end user registration
case, we show that substantial structural defects can be detected and resolved
before code generation, and that specification closure arises as a state of the
entire multi-level loop rather than of any single layer.  The paper is
intentionally framed as conceptual and architectural: the formal core is defined
here, while the reference implementation and empirical validation are placed on
the immediate roadmap.
\end{abstract}

\medskip
\noindent\textbf{Keywords:}
controlled generative development;
contract-first development;
AI-native software design;
typed dependency graphs;
specification closure;
verification-driven dialogue;
task-driven context assembly

\bigskip

% ══════════════════════════════════════════
\section{Introduction}\label{sec:intro}

Generative AI has sharply reduced the cost of writing code, but it has not
solved the more fundamental problem: how to preserve the causal link between
the author's original intent, the structure of the system, and the generated
implementation.  In most practical scenarios, an LLM receives either a
free-form prompt or weakly structured context and then produces code that is
verified only after the fact: by tests, manual review, or operation in
production.  This means that structural defects are detected too late---after
they have already been embodied in code, dependencies, and data~\cite{dakhel2023copilot,jiang2024survey}.

In this paper we propose \textbf{Controlled Generative Development (CGD)}, an
approach in which the specification precedes code, guides generation, and
outlives code as the canonical artifact.  The central object of CGD is the
formal triple

\[
  T = \langle Y,\; G,\; \Sigma \rangle,
\]
where $Y$ is a YASF specification in controlled natural language, $G$ is a
typed directed graph of functions and tables, and $\Sigma$ is a set of table
signatures.  The triple is treated not as documentation but as a compilable
and verifiable project invariant: changes in one component either cascade
consistently into the others or are flagged by the verifier as a mismatch.

CGD does not claim the status of a main-track empirical paper with already
demonstrated industrial effectiveness.  This is a conceptual and architectural
paper with a formal core, explicit limitations, and a roadmap toward a
reference implementation.  We deliberately do not claim strong operational
metrics here.  Instead, the paper answers an earlier and more foundational
question: what should the minimal formal layer look like if AI-assisted
development is to become structurally controllable before code is generated?

The main contributions of the paper are as follows.
\begin{enumerate}
  \item We introduce CGD as a distinct \textbf{contract-first} framework
        situated between prompt-first, workflow-first, spec-first, and
        formal-first approaches.
  \item We define YASF as a minimal authoring format for specification and,
        at the same time, as a potential intermediate representation for
        larger systems.
  \item We specify the graph model
        $G = (V, E, \mathrm{src}, \mathrm{dst}, \tau_V, \tau_E)$ with two
        node types, five edge types, and the six-component function contract
        $C(f) = \langle C_{in}, C_{out}, R, W, T_{in}, T_{out} \rangle$.
  \item We formalize subgraph decomposition and aggregation with a computable
        correctness condition via \textsc{ComputeSubgraphContract} and
        \texttt{refinement\_map}.
  \item We define \textbf{task-driven context assembly} as a deterministic
        function of the contractual neighborhood of the target node.
  \item We close the cycle through \textbf{contract-driven generation} and
        \textbf{reverse verification}, while explicitly distinguishing what
        is already formalized from what remains on the roadmap.
  \item We demonstrate the pre-generation part of the cycle on a single
        end-to-end user registration case in five examples.
  \item We introduce \textbf{verification-driven dialogue} as an
        infrastructural loop in which the verifier not only finds defects
        but also generates targeted questions for the user until
        specification closure is reached.
\end{enumerate}


% ══════════════════════════════════════════
\section{Existing Approaches to AI-Assisted Development}\label{sec:related}

\subsection{Classification of approaches}\label{sec:classification}

Existing approaches to AI-assisted development can be usefully classified
along two axes: the degree of specification formalization and the point at
which verification takes place.  For the purposes of this paper we distinguish
\textbf{five} categories.

\paragraph{Prompt-first.}
The developer formulates the task in free form and the LLM generates code.
There is no independent specification artifact; the prompt is ephemeral and
poorly reproducible.  Verification is entirely
post-generative~\cite{dakhel2023copilot,cursor2024}.

\paragraph{Workflow-first.}
Generation is organized as a sequence of roles, agents, or stages.  The
process becomes more structured, but the data, dependencies, and contracts are
still passed as text between steps.  Verification is still aimed at the result
rather than at a formal model of the
system~\cite{qian2024chatdev,hong2024metagpt,wu2023autogen}.

\paragraph{Spec-first.}
The specification becomes the primary artifact: the developer maintains a
structured text from which code and tests are generated.  This is a major step
forward because intent is separated from implementation.  However, the
architecture usually remains single-level: ``specification $\to$ generation
$\to$ tests $\to$ revise the specification.''  The closest industrial example
to our work is \textbf{CodeSpeak}: textual \texttt{.cs.md} specs, code
generation, mixed mode, and takeover for legacy code.  This class of tools
confirms that there is already industrial demand for specification-first
practice~\cite{breslav2026codespeak,habr2026codespeak,microsoft2023typechat,orosz2026kotlin}.

\paragraph{Formal-first.}
The system is described in a formal language with mathematical
semantics---TLA+, Alloy, Lean, Z, VDM, and similar systems.  The guarantees
are maximal, but the cost of entry is high, and direct integration with
LLM-based code generation is usually not the primary
goal~\cite{lamport2002tla,jackson2012alloy}.

\paragraph{Contract-first (CGD).}
The developer writes in the minimal YASF format, which is deterministically
mapped into the triple $T = \langle Y, G, \Sigma \rangle$.  Structural and
contractual verification is performed \textbf{before} code generation.  The
LLM context is assembled from the formal neighborhood of a node rather than
from a global bag of text.  The architecture is multi-level: different classes
of defects are filtered at different infrastructure levels.

\subsection{Evolutionary axis}\label{sec:evolution}

These categories form a working evolutionary axis in terms of increasing
formalization and depth of control:

\medskip
\centerline{\textbf{Prompt-first $\to$ Workflow-first $\to$ Spec-first $\to$
Contract-first (CGD).}}
\medskip

\noindent
At each step a new stable layer appears.  Workflow-first adds process
structure.  Spec-first introduces the specification as an explicit artifact.
Contract-first adds a formal system model and pre-generative verification.
Formal-first sits at the same end of the spectrum in terms of rigor, but
forms a parallel branch: it is not the next evolutionary step of LLM-native
development, but a separate tradition with a different balance between
guarantees and entry
cost~\cite{qian2024chatdev,hong2024metagpt,wu2023autogen,lamport2002tla,jackson2012alloy,meyer1997oosc,evans2003ddd}.

\subsection{CodeSpeak as the closest industrial analogue}\label{sec:codespeak}

CodeSpeak matters for this paper not as a competitor, but as evidence of
engineering demand.  It shows that a spec-first cycle based on \texttt{.cs.md}
specifications, generation, mixed mode, and takeover is already perceived as a
practical alternative to direct code editing.  CGD, however, operates one level
deeper: it introduces a formal multi-level scaffold in which the specification
is not only an input to the generator, but also an object of structural
verification before code
exists~\cite{breslav2026codespeak,habr2026codespeak,orosz2026kotlin}.

In CGD terms, CodeSpeak primarily covers the authoring-specification level and
part of the generation level, but it does not explicitly introduce a separate
graph layer, a separate contract-and-signature layer, or verification-driven
dialogue as an infrastructural loop.  CGD is aimed precisely at this vacant
niche~\cite{breslav2026codespeak,habr2026codespeak,orosz2026kotlin}.


% ══════════════════════════════════════════
\section{Controlled Generative Development}\label{sec:cgd}

\subsection{The triple as an invariant}\label{sec:triple}

The central artifact of CGD is the triple $T = \langle Y, G, \Sigma \rangle$,
where $Y$ captures the semantic dimension of the specification, $G$ its
structural dimension, and $\Sigma$ its typological dimension.  These three
dimensions are not strictly orthogonal: $Y$ compiles into $G$, the graph
refers to $\Sigma$, and $\Sigma$ constrains the admissible links in $G$.
Nevertheless, they form three distinguishable coordinate axes of description,
sufficiently independent to serve as a working basis for design.

\begin{thesis}[Sufficiency of the triple]\label{thesis:triple}
For the class of information systems with discrete data flows, the triple
$\langle Y, G, \Sigma \rangle$ is sufficient to express and verify those
aspects of a system that concern the structural organization of functions,
data flows, and contractual compatibility.  Full formal and empirical
validation of this hypothesis remains future work.
\end{thesis}

The determinism of compilation from YASF into graph and signatures is likewise
stated as an \textbf{implementation invariant}, to be checked by conformance
tests, rather than as an already established empirical fact.

\subsection{Task-driven context assembly}\label{sec:context}

The traditional way of working with LLMs looks like ``collect as much context
as possible, compress it, and hope the model uses exactly what matters.''
This strategy systematically creates two extremes: \emph{context pollution}
and \emph{context loss}~\cite{lewis2020rag}.

CGD replaces this pattern with \textbf{task-driven context assembly}.  For a
target function $f$, the context is defined not by dialogue history and not by
an arbitrary RAG set, but by the formal neighborhood of its contract.  In the
basic case, the context includes:
\begin{itemize}
  \item the full contract $C(f)$;
  \item the signatures of all tables in $C_{in}(f)$, $C_{out}(f)$, $R(f)$,
        $W(f)$;
  \item the contracts of first-order functions connected to $f$ through
        $T_{in}(f)$ and $T_{out}(f)$;
  \item when necessary, the signatures of tables jointly used with first-order
        neighbors.
\end{itemize}

The context therefore becomes a deterministic function of the node and its
contractual neighborhood.  It is local, reproducible, and independent of how
exactly the developer or agent happened to describe the project in earlier
interactions.

\subsection{Contract-driven generation, reverse verification, and the sterile
zone}\label{sec:sterile}

Once specification closure has been reached, the code generator receives not a
free textual description but a strict technical assignment derived from the
contract and signatures.  The contract fixes inputs, outputs, reads, writes,
and triggers.  The table signatures fix fields, keys, kind, and constraints.
This constrains the space of admissible implementations before the first token
of code is generated.

Generation itself remains non-deterministic: the model may choose different
algorithms, styles, and code structures.  But any implementation is then
compared not with the original prompt, but with the canonical triple.

This is also where the principle of the \textbf{sterile zone} is introduced.
In CGD a project is considered structurally sterile as long as unambiguous
traceability and reverse verifiability are preserved between $Y$, $G$,
$\Sigma$, and code.  Manual edits to generated code are treated as violations
of sterility.  Classical round-trip engineering is deliberately rejected as
the default process model.  If legacy or handwritten code must be included in
the system, it goes through an \textbf{assimilation} procedure:
decomposition, extraction of the actual triple $T'$, reconstruction of YASF,
and renewed achievement of closure~\cite{breslav2026codespeak}.

\subsection{Verification-driven dialogue}\label{sec:dialogue}

Verification-driven dialogue turns specification building from a monologue by
the author into a controlled cycle:

\medskip
\centerline{\textbf{build the triple $\to$ verify $\to$ classify defects
$\to$ ask the user $\to$ refine the triple $\to$ repeat.}}
\medskip

\noindent
The verifier splits defects into three classes.
\begin{itemize}
  \item \textbf{Blocking} --- defects that cannot be resolved correctly
        without a user decision.
  \item \textbf{Auto-resolvable} --- only safe normalizations that do not
        alter domain semantics: canonicalization of names, aliases,
        enum/PK/FK, kind, sections, and duplicate edges.
  \item \textbf{Advisory} --- warnings that may require a human decision but
        do not block progress of the cycle itself.
\end{itemize}

Specification closure is therefore not a property of a single level and not a
synonym for ``the graph passed L2.''  It is a state of the whole loop,
reached after the blocking queue is empty and all user decisions have cascaded
downward through the specification.

\subsection{Multi-level CGD infrastructure}\label{sec:levels}

\subsubsection{Strategic principle}\label{sec:strategy}

The strategic principle of CGD is simple: \textbf{make the infrastructure more
complex, not the author}.  The author should work with a minimal and naturally
readable representation.  Everything related to formal rigor, contractual
algebra, graph checking, and context assembly is shifted into the
infrastructural layer.

\subsubsection{Five processing levels}\label{sec:five-levels}

CGD uses five levels.

\paragraph{L0 --- user / interface.}
At this level there is only intent, expressed in free natural language.

\paragraph{L1 --- authoring / YASF.}
At this level a partial triple $T_1 = \langle Y, -, - \rangle$ appears: a
structured YASF block with a canonical subset of sections.

\paragraph{L2 --- structural / graph.}
Here $T_2 = \langle Y, G, - \rangle$ appears: YASF compiles into a typed
graph.  Structural defects are filtered out---illegal edges, isolated nodes,
duplicates, broken flows, and missing mandatory paths.

\paragraph{L3 --- contractual / signatures.}
At this level the full triple $T_3 = \langle Y, G, \Sigma \rangle$ is
formed.  Signature compatibility, PK/FK integrity, kind, enum compatibility,
and contract completeness are checked.

\paragraph{L4 --- generation + reverse verification.}
The canonical triple is used for generation and then for reverse checking of
the actual triple $T'$ extracted from code.

\subsubsection{Layered defect filtering and specification
closure}\label{sec:closure}

Each level filters its own class of defects and thereby reduces the burden on
the next one.  A problem caught at L2 should not survive until L4.  This is
why closure is formulated as a state of the whole infrastructural loop.

In this paper \textbf{specification closure} is considered achieved when the
following five conditions hold simultaneously:

\begin{enumerate}
  \item all structural checks are green;
  \item there are no unresolved references, signatures, or aliases;
  \item all contracts are closed;
  \item all mandatory paths terminate in an external effect;
  \item the blocking-defect queue is empty.
\end{enumerate}

\subsubsection{Event routing: downward cascade and upward
escalation}\label{sec:cascade}

CGD uses two complementary trajectories.

\textbf{Downward cascade}: L0 $\to$ L1 $\to$ L2 $\to$ L3 $\to$ L4.  Intent
is progressively formalized.

\textbf{Upward escalation}: if L2 or L3 detects a blocking defect, it is
raised to L0 as a targeted question to the user; the answer then cascades
back down through a YASF update, graph recompilation, and contract
recomputation.

\subsubsection{Comparison with single-level
approaches}\label{sec:single-level}

Single-level spec-first tools mix semantic incompleteness, structural defects,
contractual mismatches, and generator errors within one
loop~\cite{breslav2026codespeak}.  CGD separates these causes across
responsibility levels.  As a result, every defect receives an address: where
it arose, where it must be resolved, and how the consequences of the
resolution cascade.

\subsubsection{The triple as a working basis of the specification
space}\label{sec:basis}

The triple $\langle Y, G, \Sigma \rangle$ defines three distinguishable
dimensions of the specification space:
\begin{itemize}
  \item $Y$ covers intent, processing rules, and branching conditions;
  \item $G$ covers the topology of functions and data flows;
  \item $\Sigma$ covers table typing and integrity constraints.
\end{itemize}

We treat this triple as a \textbf{working basis} and a constructive design
thesis~\cite{meyer1997oosc,chen2023kgse}.  This makes it possible to
interpret the CGD levels as different projections of the same formal space and
makes the architecture extensible: a new level can be added without revising
the basis, provided that a mapping from the space $T$ is defined for
it~\cite{lamport2002tla,jackson2012alloy}.

The extensibility principle has a practical integration consequence.  Any
external tool---a code generator, linter, CI/CD pipeline, documentation
generator, or cross-project analyzer---can be integrated into CGD as an
additional level, provided that it defines its own mapping from the space $T$.
The triple $T = \langle Y, G, \Sigma \rangle$ serves as a common formal
integration interface: tools connect through a formal projection contract
rather than through an ad hoc API.  In this sense, CGD is an integration
platform, not an alternative to existing development tools.

\subsubsection{Representation of control
structures}\label{sec:control-flow}

Control flow is covered by the triple \emph{representationally}, not in full
operational semantics.

\begin{itemize}
  \item In $Y$, branching, loops, and parallelism are described in the
        Processing section in controlled natural language.
  \item In $G$, these structures become topology: success-path and
        error-path, back-edges for loops, and independent branches for
        parallelism.
  \item In $\Sigma$, the data traveling along these paths are typed.
\end{itemize}

What remains outside the current formalization are the exact
branch-selection rules, execution policies for loops, priorities of parallel
branches, timing, and transactionality.  This is deliberately deferred to
the limitations section.


% ══════════════════════════════════════════
\section{Formal Core}\label{sec:formal}

\subsection{YASF}\label{sec:yasf}

\textbf{YASF (Yet Another Specification Format)} is the minimal authoring
layer of specification in controlled natural language.  For the purposes of
this paper, we use a \textbf{canonical subset of sections} sufficient to
derive the graph and contracts: purpose, input data, reference data,
processing, output data, errors, side effects, triggers, and contract.

It is important not to overstate the status of the format.  YASF is not
presented here as a forever fixed set of sections.  For larger systems it may
serve not only as something written by a human, but also as an internal
\textbf{IR} generated by tooling.  What remains invariant is not the origin
of YASF, but its role as the semantic input to the triple.

Compilation $\text{YASF} \to (G, \Sigma)$ is treated as a deterministic
implementation goal and should be checked by conformance tests.

\subsection{Graph model and contracts}\label{sec:graph}

The graph is defined as

\[
  G = (V,\; E,\; \mathrm{src},\; \mathrm{dst},\; \tau_V,\; \tau_E),
\]
where $\tau_V(v) \in \{\texttt{function},\; \texttt{table}\}$, and
$\tau_E(e) \in \{\texttt{consumes},\; \texttt{produces},\;
\texttt{reads},\; \texttt{writes},\; \texttt{triggers}\}$.

In the current synchronized version of the paper, edge types are defined as
follows:
\begin{itemize}
  \item \texttt{consumes}: table $\to$ function
  \item \texttt{produces}: function $\to$ table
  \item \texttt{reads}: table $\to$ function
  \item \texttt{writes}: function $\to$ table
  \item \texttt{triggers}: function $\to$ function
\end{itemize}

Table $\to$ table links are forbidden.  Function $\to$ function links are
permitted only for \texttt{triggers}.  Triggers to tables are forbidden.

For each function $f$, the contract is defined as

\[
  C(f) = \langle C_{in},\; C_{out},\; R,\; W,\; T_{in},\; T_{out} \rangle,
\]
where all components are derived directly from the
graph~\cite{meyer1997oosc}.

\subsubsection{Subgraph and boundary}\label{sec:subgraph}

Let $V_{\text{func}}(S) \subseteq V$ be a chosen set of functions.  Then
define the set of subgraph tables as

\[
  V_{\text{tables}}(S) = \bigl\{ t \in V \;\big|\;
  \tau_V(t)=\texttt{table} \text{ and } t \text{ is incident to at least one
  edge connected to a function from } V_{\text{func}}(S) \bigr\}.
\]

Then:
\begin{itemize}
  \item $V(S) = V_{\text{func}}(S) \cup V_{\text{tables}}(S)$;
  \item $E(S) = \bigl\{ e \in E \;\big|\;
        \mathrm{src}(e) \in V(S) \text{ and } \mathrm{dst}(e) \in V(S)
        \bigr\}$.
\end{itemize}

A table is considered internal if all its incident edges connect it only to
functions from $V_{\text{func}}(S)$.  The remaining tables form the boundary
$\partial S$.

\subsubsection{\textsc{ComputeSubgraphContract}}\label{sec:compute}

\begin{lstlisting}
function ComputeSubgraphContract(S):
    Vfunc = V_func(S)
    Vtabs = V_tables(S)

    Internal = { t in Vtabs |
                 every edge incident to t
                 is incident only to functions in Vfunc }
    Boundary = Vtabs - Internal

    Cin(S)  = { t in Boundary |
                exists consumes edge  t -> f, f in Vfunc }
    Cout(S) = { t in Boundary |
                exists produces edge f -> t, f in Vfunc }
    R(S)    = { t in Boundary |
                exists reads edge    t -> f, f in Vfunc }
    W(S)    = { t in Boundary |
                exists writes edge   f -> t, f in Vfunc }
    Tin(S)  = { g not in Vfunc |
                exists trigger g -> f, f in Vfunc }
    Tout(S) = { g not in Vfunc |
                exists trigger f -> g, f in Vfunc }

    return <Cin(S), Cout(S), R(S), W(S), Tin(S), Tout(S)>
\end{lstlisting}

This definition is consistent with the examples: external tables are present
in $V(S)$ as boundary tables, and the subgraph contract is computed precisely
from them rather than from abstract ``interface tables'' outside the induced
subgraph.

\subsubsection{Decomposition and aggregation}\label{sec:decomp}

A function $f$ may be replaced with a subgraph $S$ if three conditions hold:
\begin{enumerate}
  \item $S$ is induced by its $V(S)$;
  \item all external links pass only through $\partial S$ and the trigger
        boundary;
  \item the external contract, after applying the
        internalization/externalization rule, is preserved: $C(S) = C(f)$.
\end{enumerate}

One clarification is needed here.  If a function $g \in T_{out}(f)$ is
included inside the subgraph, \textbf{trigger internalization} occurs: $g$ is
removed from $T_{out}(S)$, and its boundary outputs become part of
$C_{out}(S)$.  Aggregation performs the inverse
operation---\textbf{externalization}.  Thus the reversibility of
decompose/aggregate becomes parameterized by the subgraph boundary, but is
not broken.

Aggregation is admissible if the subgraph is collapsed into a node $f_S$ with
the same external contract and is accompanied by a \texttt{refinement\_map}
sufficient to recover the internal structure exactly.

\subsection{Table signatures and compatibility}\label{sec:signatures}

For each table $t$, the signature is defined as

\[
  \Sigma(t) = \langle \text{Fields},\; \text{Keys},\;
  \text{Kind},\; \text{Constraints} \rangle.
\]

In the current JSON profile, fields are typed by base types such as
\texttt{string}, \texttt{integer}, \texttt{boolean}, \texttt{float},
\texttt{timestamp}, \texttt{date}, \texttt{uuid}, \texttt{text}, and
\texttt{json}; enum and regex constraints are expressed via
\texttt{Constraints}.  The attribute \texttt{role} may be present as an
auxiliary semantic hint, but in the current version it \textbf{does not
participate in formal compatibility}.

The canonical kinds in the synchronized version are \textbf{entity, event,
reference, log, projection, error}.  For the purposes of the current version,
every table, including \texttt{projection}, is required to have a canonical PK
or an explicit surrogate key; this removes the contradiction between
representation semantics and serialization rules.

Signature compatibility is checked after canonicalization.  It is convenient
to distinguish \textbf{strict equality} $\Sigma(A) = \Sigma(B)$ from
\textbf{inclusion compatibility} $\Sigma(A) \supseteq \Sigma(B)$.  It is
important to state the direction of the enum rule explicitly.  If a producer
function emits a table that is then interpreted by a consumer function, then
for every enum field the following must hold:

\[
  \text{Allowed}(\text{producer}) \subseteq \text{Allowed}(\text{consumer}).
\]
Otherwise the producer may generate a value that the consumer cannot
interpret, which is an L3 defect.

In the current version of the paper, per-consumer field contracts are
\textbf{not formalized}: a function operates on a table as a whole object,
rather than on a formally specified projection of its fields.


% ══════════════════════════════════════════
\section{End-to-End Case: User Registration}\label{sec:case}

\subsection{Example 1 --- Birth of the triple}\label{sec:ex1}

From the request ``register a user by email, persist the record, send a
confirmation, and log the event,'' the following base tables are extracted:
\textbf{RegistrationRequest}, \textbf{User}, \textbf{RegistrationEvent},
\textbf{AuditLog}, \textbf{NotificationEvent}; and two functions:
\textbf{RegisterUser} and \textbf{SendConfirmation}.

Top-level contract:

\[
  C(\text{RegisterUser}) = \langle
  \{\text{RegistrationRequest}\},\;
  \{\text{User},\, \text{RegistrationEvent}\},\;
  \{\text{User}\},\;
  \{\text{AuditLog}\},\;
  \emptyset,\;
  \{\text{SendConfirmation}\}
  \rangle.
\]

Trigger-function contract:

\[
  C(\text{SendConfirmation}) = \langle
  \{\text{User}\},\;
  \{\text{NotificationEvent}\},\;
  \emptyset,\;
  \emptyset,\;
  \{\text{RegisterUser}\},\;
  \emptyset
  \rangle.
\]

The graph contains \textbf{2~functions, 5~tables, and 8~edges}.

\subsection{Example 2 --- Decomposition}\label{sec:ex2}

The function \textbf{RegisterUser} is decomposed into
\textbf{ValidateInput}, \textbf{CheckUniqueness}, \textbf{CreateUser}, and
\textbf{SendConfirmation}.  Two intermediate tables appear inside the
subgraph: \textbf{ValidatedRequest} and \textbf{UniqueRequest}.

Applying \textsc{ComputeSubgraphContract} yields:

\[
  C(S) = \langle
  \{\text{RegistrationRequest}\},\;
  \{\text{User},\, \text{RegistrationEvent},\, \text{NotificationEvent}\},\;
  \{\text{User}\},\;
  \{\text{AuditLog}\},\;
  \emptyset,\;
  \emptyset
  \rangle.
\]

This is consistent with the internalization rule:
\textbf{SendConfirmation} has moved inside the subgraph, so it disappears
from $T_{out}(S)$, while its boundary output \textbf{NotificationEvent}
appears in $C_{out}(S)$.

At this point closure is still not reached: the success path is formalized,
but the error path is missing.  The subgraph contains
\textbf{4~functions, 7~tables, and 12~edges}.

\subsection{Example 3 --- Defect detection and repair}\label{sec:ex3}

At this step the verifier finds \textbf{2~blocking and 2~advisory} defects.

\textbf{Blocking~1.}  \textbf{ValidateInput} has no error path.

\textbf{Blocking~2.}  \textbf{CheckUniqueness} has no error path.

After two targeted questions, the user chooses a unified behavior: in both
cases the function produces \textbf{RegistrationError}.

\textbf{Advisory~1.}  There is no write into \textbf{AuditLog} on the error
path.  This is not an auto-resolvable case: the decision depends on the
business audit policy.  In the running case, the user decides \textbf{to add}
AuditLog writes for the error path.

\textbf{Advisory~2.}  \textbf{RegistrationEvent} has no consumer inside the
current graph.  The user explicitly states that the event is consumed by
external analytics.

After the update, the subfunction contracts become closed, and the top-level
contract is refined to version~\textbf{v2}:

\[
  C(\text{RegisterUser}) = \langle
  \{\text{RegistrationRequest}\},\;
  \{\text{User},\, \text{RegistrationEvent},\, \text{RegistrationError}\},\;
  \{\text{User}\},\;
  \{\text{AuditLog}\},\;
  \emptyset,\;
  \{\text{SendConfirmation}\}
  \rangle.
\]

Specification closure is reached here, not at the decomposition step.  The
subgraph now contains \textbf{4~functions, 8~tables, and 16~edges}.

\subsection{Example 4 --- Aggregation}\label{sec:ex4}

The repaired subgraph is collapsed back into the single node
\textbf{RegisterUser}.  During aggregation, the inverse trigger
externalization occurs: \textbf{SendConfirmation} once again becomes an
external function in $T_{out}(\text{RegisterUser})$, and
\textbf{NotificationEvent} moves back into \textbf{SendConfirmation}'s
contract.

The external contract of the aggregated node matches~\textbf{v2}.  The
\texttt{refinement\_map} preserves the internal functions, intermediate
tables, and internal edges.

It is important here to distinguish two levels of reversibility.  At the
topology level, the transformation decompose/aggregate is reversible.  At
the level of contract component distribution, reversibility is
\textbf{parameterized by the subgraph boundary}: an internalized trigger
moves its outputs into $C_{out}(S)$, and an externalized trigger returns
them to the child function.  This is not a side effect of an informal
heuristic, but part of the transformation rule itself.

After aggregation, the graph contains \textbf{2~functions, 6~tables, and
9~edges}, while closure is preserved.

\subsection{Example 5 --- Model comparison, normalization, and convergence
to the project canon}\label{sec:ex5}

The same request is sent to two models.
\begin{itemize}
  \item \textbf{Model~A} returns a simplified graph with names such as
        \texttt{InputForm}, \texttt{UserRecord}, and \texttt{EmailLog},
        without PK/kind and without an error path.
  \item \textbf{Model~B} returns a more detailed structure with
        \texttt{RegistrationRequest}, \texttt{Users},
        \texttt{Notification}, and \texttt{AuditLog}, but still without
        \texttt{RegistrationError} and \texttt{RegistrationEvent}.
\end{itemize}

Comparison is performed in three layers.
\begin{enumerate}
  \item \textbf{Overall topology.}  Both models reproduce the success path
        ``input $\to$ processing $\to$ user creation $\to$ notification.''
  \item \textbf{Structural differences.}  Names, signature completeness, and
        the presence of reads/writes differ.
  \item \textbf{Deviations from the project canon.}  Both models lack the
        error path and \texttt{RegistrationEvent}; Model~A is additionally
        structurally incomplete.
\end{enumerate}

Normalization does not mean that the models ``converge by themselves.''  On
the contrary, triples are extracted from both outputs, discrepancies are
classified, and then both are normalized against the \textbf{project
reference model} from Examples~1--4.  The canonical result matches the
\textbf{v2}~contract.

After normalization, the graph has \textbf{2~functions, 6~tables, and
9~edges}.

\subsection{Quantitative summary}\label{sec:quant}

All numbers in this section are \textbf{illustrative}, not benchmark data.

\begin{table}[h]
\centering
\begin{tabular}{@{}lccc@{}}
\toprule
\textbf{Example} & \textbf{Functions} & \textbf{Tables} & \textbf{Edges} \\
\midrule
1 --- Birth        & 2 & 5 & 8  \\
2 --- Decomposition & 4 & 7 & 12 \\
3 --- Defect repair & 4 & 8 & 16 \\
4 --- Aggregation   & 2 & 6 & 9  \\
5 --- Normalization & 2 & 6 & 9  \\
\bottomrule
\end{tabular}
\caption{Graph metrics across the five running examples.}
\label{tab:metrics}
\end{table}

In terms of defects, the case yields \textbf{2~blocking + 2~advisory}
defects and one iteration of verification-driven dialogue before closure is
reached.  For the function \textbf{CheckUniqueness}, task-driven context
assembly reduces local context by approximately \textbf{65\%} relative to the
full graph; this number is illustrative and is not presented as a general
empirical metric.


% ══════════════════════════════════════════
\section{Verification Without Code}\label{sec:verification}

\subsection{Pre-generative verification}\label{sec:pregen}

Before code generation, the following are checked:
\begin{itemize}
  \item type correctness of edges;
  \item uniqueness of identifiers and absence of duplicates;
  \item absence of isolated nodes;
  \item presence of PK and kind for every table;
  \item consistency of contracts with the graph;
  \item completeness of mandatory paths;
  \item canonicality of signatures and contractual compatibility.
\end{itemize}

The key effect of CGD is that a significant part of architectural defects
becomes observable \textbf{before} code exists.

\subsection{Post-generative verification}\label{sec:postgen}

The post-generative part of the cycle is specified in this paper as a formal
roadmap rather than as a completed reference implementation.  The actual
triple $T'$ is to be extracted from code, after which $T'$ is compared with
the canonical~$T$.

Some such discrepancies can be extracted strictly: extra dependencies, missing
outputs, extra writes, altered signatures.  Some remain analytical: hidden
logic, non-formalized path assumptions, or unused fields in the absence of
per-consumer field contracts.


% ══════════════════════════════════════════
\section{Limitations and Roadmap}\label{sec:limits}

\subsection{Current limitations}\label{sec:current-limits}

The current version does not
support~\cite{lamport2002tla,evans2003ddd,chen2023kgse}:
\begin{itemize}
  \item per-consumer field contracts;
  \item path assumptions as a formal contract component;
  \item contract subtyping;
  \item developed operational semantics for loops and parallelism;
  \item concurrent access and temporal annotations;
  \item schema migrations and a full-fledged runtime-monitoring layer.
\end{itemize}

\subsection{Implementation status}\label{sec:impl-status}

At the time of this paper, YASF, graph~v0.3, tables~v0.2, JSON
representations, and the running case are formally defined.  A reference
implementation of the validator, compiler, and reverse verification is not yet
claimed as a production-ready system.  This boundary is drawn consciously and
explicitly.

\subsection{Near-term steps}\label{sec:roadmap}

The roadmap includes:
\begin{itemize}
  \item a reference implementation of the YASF $\to$ JSON-artifact compiler;
  \item a structural + contract consistency validator;
  \item a task-driven context generator;
  \item a reverse verification prototype;
  \item case expansion beyond the user-registration scenario;
  \item controlled evaluation on several real projects.
\end{itemize}

\subsection{Threats to validity}\label{sec:threats}

\textbf{Construct validity.}  The notions of structural defect, closure, and
contractual compatibility are introduced by the authors and may fail to cover
all practical error classes.

\textbf{Internal validity.}  The running case is synthetic and small.  It
tests internal consistency of the model, not industrial performance.

\textbf{External validity.}  Transferability to domains with complex
concurrency, temporal logic, and rich state semantics has not yet been
demonstrated.

\textbf{Reliability.}  The quality of verification-driven dialogue depends
not only on the formal model, but also on how well the infrastructure turns a
defect into a clear and targeted question for the
user~\cite{evans2003ddd,jiang2024survey}.


% ══════════════════════════════════════════
\section{Conclusion}\label{sec:conclusion}

CGD offers not just another prompt template and not just another workflow
around an LLM, but an earlier control layer---a layer at which the
specification exists as a formal and verifiable invariant even before code is
written~\cite{jackson2012alloy}.

The central result of the paper is the construction
$T = \langle Y, G, \Sigma \rangle$ together with the associated multi-level
infrastructure L0--L4.  This infrastructure makes it possible to filter
defects layer by layer, assemble task-oriented context from the contractual
neighborhood, keep the project inside the sterile zone, and avoid delegating
to the code generator those architectural decisions that should have been made
earlier.  The architecture is open to extension: new levels and external tools
can be connected through a formal projection contract, turning CGD into an
integration platform.

We deliberately state the completeness of the triple and the sufficiency of
its basis as a \textbf{design thesis}, not as an already proven mathematical
fact.  The immediate next step is a reference implementation and controlled
empirical validation.  Even in its current form, however, CGD defines a clean
architectural framework: the specification should precede code, govern
generation, and outlive code as the canonical artifact.


% ── Bibliography ──
\bigskip
\noindent\textbf{Supplementary materials.}
Appendices, JSON schemas, and examples are available at
\url{https://github.com/alexfdo/cgd-paper}.


\bibliographystyle{plainnat}
\bibliography{references}

\end{document}
